%! Orthogonality of the Four Subspaces — Unit 4, Lesson 4.1 — Guided Notes (Answer Key)
\documentclass[10pt]{article}
\usepackage{linalg-article}
\usepackage{linalg-key}   % pulls in -boxes; do NOT also load -boxes
\newcommand{\TallMath}[1]{$\displaystyle #1\rule[-1.4em]{0pt}{3.2em}$}
\newcommand{\vv}{\mathbf{v}}
\newcommand{\ww}{\mathbf{w}}
\newcommand{\xx}{\mathbf{x}}
\newcommand{\yy}{\mathbf{y}}
\newcommand{\bb}{\mathbf{b}}
\newcommand{\zero}{\mathbf{0}}
\newcommand{\T}{\mathsf{T}}

\begin{document}

\pageheader{Unit 4, Lesson 4.1}{Guided Notes --- Answer Key}
\namedateperiod

\vspace{0.12in}

\begin{objectivebox}
\textbf{By the end of this lesson, I will be able to\dots}
\begin{itemize}\setlength{\itemsep}{2pt}
  \item say what it means for two \emph{subspaces} to be \emph{orthogonal}, and test it;
  \item explain \emph{why} the row space $\perp$ the nullspace and the column space $\perp$ the left nullspace;
  \item recognize the two pairs as \emph{orthogonal complements} --- perpendicular \emph{and} filling the whole space.
\end{itemize}
\end{objectivebox}

\vspace{0.06in}

\begin{vocabbox}
\small
Fill in each term as we build it together.
\vspace{2pt}
\vocabans{Orthogonal subspaces}{subspaces $V,W$ where every vector in $V$ is perpendicular to every vector in $W$ ($\vv\cdot\ww=0$); enough to check basis against basis}
\vocabans{Orthogonal complement $V^{\perp}$}{\emph{all} vectors perpendicular to $V$; perpendicular to $V$ \emph{and} dimensions fill the space, so they meet only at $\zero$}
\vocabans{The two perpendicular pairs (§3.6, now in full)}{$N(A)\perp C(A^{\T})$ in $\mathbb{R}^n$ and $N(A^{\T})\perp C(A)$ in $\mathbb{R}^m$; each pair is a complement pair}
\end{vocabbox}

\begin{hookbox}
Picture a flat tabletop inside a room (a \emph{plane} in $\mathbb{R}^3$).
\begin{itemize}\setlength{\itemsep}{1pt}
  \item Which directions point at a \emph{right angle} to the \emph{entire} tabletop? \ansline{Straight up/down --- the single direction perpendicular to the whole plane.}
  \item How many independent such directions are there --- a line, a plane, \dots? \ansline{Just one --- a line (dim $1$). The plane (dim $2$) $+$ the line (dim $1$) fill $\mathbb{R}^3$.}
\end{itemize}
\end{hookbox}

\begin{notesbox}{1. From Orthogonal Vectors to Orthogonal Subspaces}
Two \emph{vectors} are perpendicular when $\vv\cdot\ww=0$ (Lesson~1.2). We now lift this to whole
subspaces. Two subspaces $V$ and $W$ are \textbf{orthogonal} when \emph{every} vector in $V$ is
perpendicular to \ans{every vector} in $W$.

\vspace{3pt}
\textbf{A shortcut.} You do not have to test infinitely many vectors --- it is enough to check that
every \textbf{basis} vector of $V$ is perpendicular to every basis vector of $W$ (dot products spread
over sums).

\vspace{3pt}
\textbf{Careful --- perpendicular-looking is not enough.} Two walls of a room \emph{meet} along a
vertical line, and a vector \emph{along} that shared line lies in both walls --- it cannot be
perpendicular to itself (unless it is $\zero$). So two walls are \ans{not} orthogonal. Orthogonal
subspaces may only share the point \ans{$\zero$}.
\end{notesbox}

\begin{notesbox}{2. Why the Row Space $\perp$ the Nullspace}
Take any $\xx$ in the nullspace, so $A\xx=\zero$. Read this equation \textbf{one row at a time}:
\[
  (\text{row }1)\cdot\xx=0,\quad (\text{row }2)\cdot\xx=0,\quad \dots,\quad (\text{row }m)\cdot\xx=0.
\]
So $\xx$ is perpendicular to \emph{every row}. And if $\xx\perp$ every row, then $\xx\perp$ every
\ans{combination} of the rows --- that is, $\xx\perp$ the whole \ans{row space}. Therefore
\[
  N(A)\ \perp\ C(A^{\T}) \qquad\text{(nullspace $\perp$ row space, inside }\mathbb{R}^n).
\]
Run the very same argument on $A^{\T}$ (its rows are the \emph{columns} of $A$) to get the second pair:
\[
  N(A^{\T})\ \perp\ C(A) \qquad\text{(left nullspace $\perp$ column space, inside }\mathbb{R}^m).
\]
This is the \emph{why} behind the §3.6 picture: orthogonality is just $\text{dot}=0$, one row at a time.
\end{notesbox}

\begin{notesbox}{3. Orthogonal Complements: Perpendicular \emph{and} Filling the Space}
Perpendicular alone is not the whole story. Inside $\mathbb{R}^n$ the row space has dimension $r$ and
the nullspace has dimension $n-r$, and
\[
  r+(n-r)=\ans{$n$}.
\]
The two perpendicular dimensions \emph{use up the entire space}. When two subspaces are perpendicular
\textbf{and} their dimensions add to fill the space (so they meet only at $\zero$), each is the
\textbf{orthogonal complement} of the other, written $V^{\perp}$. The nullspace is \emph{every} vector
perpendicular to the row space --- no more, no less. The four fundamental subspaces are therefore
\textbf{two orthogonal-complement pairs}, one in each room:

\begin{center}
\begin{tikzpicture}[scale=1.0, every node/.style={font=\footnotesize}]
  % Input room
  \draw[burgundy, thick, rounded corners] (-0.3,-0.3) rectangle (4.3,4.9);
  \node[burgundy] at (2,5.25) {\textbf{Input space } $\mathbb{R}^n$};
  \node[draw=burgundy, fill=blush, rounded corners, align=center, minimum width=3.5cm] (rs) at (2,3.7) {Row space $C(A^{\T})$\\ dim $r$};
  \node[draw=burgundy, fill=blush, rounded corners, align=center, minimum width=3.5cm] (ns) at (2,1.0) {Nullspace $N(A)$\\ dim $n-r$};
  \node[text=burgundy] at (2,2.35) {$\perp$};
  % Output room
  \draw[royalblue, thick, rounded corners] (6.7,-0.3) rectangle (11.3,4.9);
  \node[royalblue] at (9,5.25) {\textbf{Output space } $\mathbb{R}^m$};
  \node[draw=royalblue, fill=skyblue, rounded corners, align=center, minimum width=3.5cm] (cs) at (9,3.7) {Column space $C(A)$\\ dim $r$};
  \node[draw=royalblue, fill=skyblue, rounded corners, align=center, minimum width=3.5cm] (lns) at (9,1.0) {Left nullspace $N(A^{\T})$\\ dim $m-r$};
  \node[text=royalblue] at (9,2.35) {$\perp$};
  % Arrow A
  \draw[->, very thick, charcoal] (4.5,2.3)--(6.5,2.3) node[midway, above]{$A$};
\end{tikzpicture}
\end{center}

\emph{Not} every perpendicular pair is a complement: two perpendicular \emph{lines} in $\mathbb{R}^3$
have dimensions $1+1=2\neq 3$, so they leave directions out --- perpendicular, but not complements.
\end{notesbox}

\begin{notesbox}{4. Worked Example --- All Four Subspaces of One Matrix}
Let $A=\begin{bsmallmatrix} 1 & 1 & 2\\ 2 & 1 & 3\\ 3 & 2 & 5 \end{bsmallmatrix}$ (our Unit~3 matrix),
which has rank $r=2$, with $m=n=3$. Elimination gave the nullspace direction
$\xx=\begin{bsmallmatrix} -1\\-1\\1 \end{bsmallmatrix}$ and the left-nullspace direction
$\yy=\begin{bsmallmatrix} 1\\1\\-1 \end{bsmallmatrix}$ (because row$_3=$ row$_1+$ row$_2$).

\vspace{2pt}
\textbf{Input room} ($\mathbb{R}^3$): row space is a \emph{plane} (dim $2$), nullspace is a
\emph{line} (dim $1$). Check row$_1\cdot\xx$:
$\begin{bsmallmatrix} 1\\1\\2 \end{bsmallmatrix}\cdot\begin{bsmallmatrix} -1\\-1\\1 \end{bsmallmatrix}
= -1-1+2=\ans{$0$}$. Dimensions: $2+1=\ans{$3$}$, so they \emph{fill} $\mathbb{R}^3$ --- a
complement pair.

\vspace{2pt}
\textbf{Output room} ($\mathbb{R}^3$): column space is a plane (dim $2$), left nullspace a line
(dim $1$). Check column$_1\cdot\yy$:
$\begin{bsmallmatrix} 1\\2\\3 \end{bsmallmatrix}\cdot\begin{bsmallmatrix} 1\\1\\-1 \end{bsmallmatrix}
= 1+2-3=\ans{$0$}$. Again $2+1=3$: perpendicular \emph{and} filling --- complements.
\end{notesbox}

\begin{practicebox}
Show your dot products.
\begin{enumerate}[leftmargin=1.6em, itemsep=4pt]
  \item For $A=\begin{bsmallmatrix} 1 & 2 & 2\\ 2 & 4 & 5 \end{bsmallmatrix}$, the nullspace direction is
        $\xx=\begin{bsmallmatrix} -2\\1\\0 \end{bsmallmatrix}$. Dot $\xx$ with \emph{each} row to confirm
        $N(A)\perp C(A^{\T})$. \ansline{Row$_1$: $(1)(-2)+(2)(1)+(2)(0)=0$; Row$_2$: $(2)(-2)+(4)(1)+(5)(0)=0$. Both $0$, so $\xx\perp$ the row space.}
  \item An $m\times n$ matrix has $m=4$, $n=6$, rank $r=3$. Give the dimension of each of the four
        subspaces, and name the two orthogonal-complement pairs. \ansline{Row space $r=3$ and nullspace $n-r=3$ (in $\mathbb{R}^6$, $3+3=6$); column space $r=3$ and left nullspace $m-r=1$ (in $\mathbb{R}^4$, $3+1=4$). Pairs: $C(A^{\T})\perp N(A)$ and $C(A)\perp N(A^{\T})$.}
  \item Two perpendicular lines sit in $\mathbb{R}^3$. Are they orthogonal complements? Explain using
        dimensions. \ansline{No --- $1+1=2\neq 3$, so they do not fill $\mathbb{R}^3$; perpendicular but not complements.}
\end{enumerate}
\end{practicebox}

\begin{teachernote}
The one new idea past §3.6 is \emph{complement}: perpendicular \textbf{plus} dimensions filling the
space. Section~2 gives the ``why perpendicular'' (read $A\xx=\zero$ row by row --- dot $=0$ again);
Section~3 adds ``why they fill'' ($r+(n-r)=n$). Watch for two errors: (a) calling any two
perpendicular subspaces ``complements'' (they must also fill the space --- the two-lines counterexample),
and (b) forgetting that a vector in the \emph{intersection} of two subspaces would have to be
perpendicular to itself, forcing it to be $\zero$. The worked matrix is the Unit~3 spine, so students
can lean on numbers they already trust and spend their attention on the orthogonality.
\end{teachernote}

\end{document}
